% Figure: the apoptotic caspase cascade. Two trigger paths (intrinsic /
% mitochondrial; extrinsic / death receptor), converging on the
% executioner caspases and their substrates.
\begin{figure}[htbp]
\centering
\begin{tikzpicture}[
  trig/.style={draw=engred!80!black, fill=engred!15, very thick,
    rectangle, minimum width=2.4cm, minimum height=0.9cm,
    align=center, font=\small, rounded corners=3pt},
  signal/.style={draw=engamber!80!black, fill=engamber!15, very thick,
    rectangle, minimum width=2.4cm, minimum height=0.9cm,
    align=center, font=\small, rounded corners=3pt},
  exec/.style={draw=engblue!80!black, fill=engblue!18, very thick,
    rectangle, minimum width=2.4cm, minimum height=0.9cm,
    align=center, font=\small, rounded corners=3pt},
  outc/.style={draw=engblack, fill=engblue!8, very thick,
    rectangle, minimum width=2.6cm, minimum height=0.9cm,
    align=center, font=\small, rounded corners=3pt},
  arr/.style={-{Latex[length=2.2mm]}, thick, draw=engblack},
  inh/.style={-{Bar[width=3mm]}, thick, draw=engred},
]

% Intrinsic (left) path
\node[trig] (dna) at (0, 4.2) {DNA damage \\ \scriptsize (p53)};
\node[signal] (bcl) at (0, 2.8) {Bcl-2 / Bax \\ balance};
\node[signal] (cyt) at (0, 1.4) {cytochrome $c$ \\ release};
\node[signal] (c9) at (0, 0.0) {caspase-9 \\ (initiator)};

% Extrinsic (right) path
\node[trig] (fas) at (5.5, 4.2) {Fas / TNF-$\alpha$ \\ \scriptsize ligation};
\node[signal] (disc) at (5.5, 2.8) {DISC \\ assembly};
\node[signal] (c8) at (5.5, 0.0) {caspase-8 \\ (initiator)};

% Convergence (centre / bottom)
\node[exec] (c3) at (2.75, -1.6) {caspase-3 / -7 \\ (executioner)};
\node[outc] (mem) at (-0.5, -3.4) {membrane \\ blebbing};
\node[outc] (dna2) at (2.75, -3.4) {DNA \\ fragmentation};
\node[outc] (eat) at (6.0, -3.4) {phagocytosis \\ (no inflammation)};

% Forward arrows: intrinsic
\draw[arr] (dna) -- (bcl);
\draw[arr] (bcl) -- (cyt);
\draw[arr] (cyt) -- (c9);
\draw[arr] (c9) -- (c3);

% Forward arrows: extrinsic
\draw[arr] (fas) -- (disc);
\draw[arr] (disc) -- (c8);
\draw[arr] (c8) -- (c3);

% Executioner -> outcomes
\draw[arr] (c3) -- (mem);
\draw[arr] (c3) -- (dna2);
\draw[arr] (c3) -- (eat);

% Inhibition: Bcl-2 inhibits cytochrome c release (negative regulator
% drawn as bar)
\node[font=\scriptsize, align=center, draw=engred!80!black,
  fill=engred!8, rounded corners=2pt] (bcl2)
  at (-2.8, 2.1) {Bcl-2 \\ (survival)};
\draw[inh] (bcl2) -- (cyt);

% Cross-talk: caspase-8 can also cleave Bid -> mitochondrial path
\draw[arr, densely dashed] (c8.west) to[bend right=20]
  node[midway, above, font=\scriptsize] {Bid} (cyt.east);

\end{tikzpicture}
\caption[The apoptotic caspase cascade]{The two principal triggers
of apoptosis and their convergence on the executioner caspases. The
\textbf{intrinsic} (mitochondrial) path begins with DNA damage
detected by p53, shifts the Bcl-2 / Bax balance, releases cytochrome
$c$ from the mitochondrial inner-membrane space, and activates
caspase-9. The \textbf{extrinsic} (death-receptor) path begins with
Fas ligand or TNF-$\alpha$ binding, assembles the death-inducing
signalling complex (DISC), and activates caspase-8. Both initiators
cleave and activate caspase-3 and caspase-7, the executioner caspases,
which then cleave a few hundred substrate proteins, fragment DNA, and
drive membrane blebbing without releasing pro-inflammatory contents
into the surrounding tissue. Bcl-2 inhibits cytochrome $c$ release
(red bar); caspase-8 also cleaves Bid, linking the extrinsic path to
the mitochondrial path (dashed) \cite{paper:v6c01-vaux-korsmeyer-1999,
paper:v6c01-galluzzi-2018}.}
\label{fig:vol06:ch01:apoptosis-cascade}
\end{figure}
